\section{Empirical Portfolio Learnability}
\label{sec:empirics}

\subsection{Data, Representation, and Investment Performance}

The empirical exercise asks how a representation of stock characteristics,
finite return history, and regularization jointly shape a portfolio policy.
We use one monthly U.S. equity panel, one chronology, and one exposure scale
for all performance and attribution results. The central comparison remains
Linear, Gaussian, and Mat\'ern-$3/2$. A single trained neural portfolio policy
provides a secondary check using a learned representation. The evidence is
descriptive historical evidence; neither a favorable Sharpe ratio nor an
estimated number of active directions certifies an implementable opportunity.

Formation observations run from January 1963 through December 2024. Starting
from the JKP dictionary, we exclude nano stocks, use the 130 characteristics
with greatest coverage in 1963--1972, and remove observations missing more than
30\% of those inputs. Observed values receive average cross-sectional ranks
on $[-0.5,0.5]$; remaining missing inputs receive zero, a neutral value that
need not equal the median in the presence of ties. We accept the retrospectively
constructed characteristic dictionary as the academic information set.

There is a separate, material sample limitation. The retained panel excludes
unavailable subsequent payoffs before ranking and calculating $N_t$. A
formation-only reconstruction leaves \MissingCount\ unidentified stock-month
payoffs out of \FormationCount\ (\MissingShare\%). Exact adjacent-calendar
source matches recover none of these outcomes. Its formation ranks are
invariant to perturbing future payoff availability, whereas the retained
complete-payoff sample is explicitly conditional on that availability. We
preserve this limitation in every result rather than inventing returns or
calling the record certified real-time trading evidence. The appendix
quantifies coverage and specifies the missing source information.

The first inner training period uses 1963--1972 formations and validation
uses 1973--1977. Each subsequent January-close decision reserves the preceding
60 formation months for validation, then refits using the full available
history. December's payoff is known at January close. A decision-year block
therefore realizes returns from February through the following January; the
common evaluation period is February 1978--January 2025, comprising 564 months.
This convention governs model selection, holdings, costs, and figure labels.

The Linear policy uses a constant and the characteristic coordinates.
Gaussian and Mat\'ern-$3/2$ use 10,000 random Fourier features from fixed seed-0
banks, with lengthscale equal to the initial-sample median distance. Their
kernel definitions impose different smoothness restrictions, but these finite
feature banks do not establish a strict nesting of implemented policies.
An observed performance ranking consequently does not identify a population
approximation gap. Figure~\ref{fig:equity_kernels} compares their wealth on
identical realized months.

\begin{figure}[H]
\centering
\includegraphics[width=\textwidth]{figures/fig01_kernel_wealth.pdf}
\caption{Three representations on a common protocol. Total wealth compounds
cash returns plus the policy's excess payoff; the right panel adds illustrative
25-basis-point commissions on both sides of drift-adjusted trades and a
30-basis-point annual short fee. Initial NAV is one. These are conditional
research-sample friction sensitivities, not observed execution costs.}
\label{fig:equity_kernels}
\end{figure}

Estimation follows the maximum-Sharpe criterion, in its quadratic formulation
given by the earlier definition of $Q$. The implementation fits response-one
ridge policies and selects the penalty minimizing validation $Q$ at fitted
scale. This operation is preserved; maximizing validation Sharpe would be a
different selection rule. Each representation's 120-candidate penalty grid
is frozen from its initial training spectrum. An explicit Linear intercept
is penalized with the other coefficients; Fourier features have no separately
added unpenalized intercept. These choices enter the managed-payoff operator
and are retained in every spectral calculation.

Table~\ref{tab:kernel_performance} reports the common-period performance and
investment scale. A common positive constant, calibrated from the initial
Linear validation policy's median gross exposure, applies to every policy.
It is fixed before the first test payoff and involves neither future volatility
nor full-sample gross exposure. This permits comparison of costs while leaving
gross Sharpe invariant to the constant. The higher nonlinear point estimates
are accompanied by substantial trading requirements; representation richness
alone gives no guarantee of net performance.

\input{tables/performance}

The neural portfolio policy with ridge readout learns two 64-unit tanh hidden
layers from the same direct monthly maximum-Sharpe criterion. All stocks in
a month are aggregated before the squared criterion is differentiated.
Earlier validation chooses between two stopping budgets and five exact-readout
penalty ratios; the selected specification is refitted before deployment,
without reselection on incorporated validation observations. Its seed-0
gross and net Sharpes are \NeuralGross\ and \NeuralNet. Frozen hidden features
define a positive-semidefinite feature kernel, and the ridge head has an exact
algebraic kernel interpretation. This is a data-dependent learned-feature
kernel, distinct from the neural tangent kernel; conditional head complexity
excludes learning the hidden layers. The policy does not inherit the paper's
fixed-kernel rate. Optimization and seed diagnostics appear in the appendix.

For Gaussian minus Linear, the gross Sharpe difference is \KernelGrossDelta\;
the paired 12-month circular-block interval is [\KernelGrossLow,\KernelGrossHigh].
With the stated frictions, the difference is \KernelNetDelta\, with interval
[\KernelNetLow,\KernelNetHigh]. These intervals condition on estimated paths,
use identical resampled months for both policies, and are not full-pipeline or
multiple-comparison-adjusted inference. They motivate cautious economic
comparison rather than a claim that one representation uniformly dominates.

\subsection{Economic Spectrum and Active Complexity}

The managed-payoff spectrum orders directions by their uncentered historical
payoff second moments. An eigenvalue measures the variation supported by a
specified representation and sample. It is not invested capital, expected
return, Sharpe, or a uniquely named anomaly. Figure~\ref{fig:spectrum_2024}
holds the historical cutoff fixed before the final evaluation year. Raw
eigenvalues preserve representation units; normalized mass provides a
separate concentration diagnostic and does not rescale the fitted policies.
Every reported spectrum uses pretest observations. Its dependence on
representation, sample length, and scale means that different mass pictures
should be read alongside the fitted penalty and the subsequent portfolio,
rather than as standalone investment rankings.

\begin{figure}[H]
\centering
\includegraphics[width=\textwidth]{figures/fig02_managed_spectrum.pdf}
\caption{Gaussian and Mat\'ern-$3/2$ managed-payoff spectra at January 2024.
The common history contains formations through December 2023 and known payoffs
through January 2024. The left panel uses raw uncentered second-moment units;
the right divides by each representation's trace. Normalization is diagnostic
only; penalties and fitted positions are unchanged.}
\label{fig:spectrum_2024}
\end{figure}

Mass concentration and active complexity are distinct. Ridge weights a
direction by $\widehat\mu_j/(\widehat\mu_j+\lambda)$, so weak eigenvalues can
contribute many effective degrees of freedom while representing little
second-moment mass. The historical average mass in the first 10\% of positive
Gaussian ranks is \LeadingMass\%, but that group contributes only
\LeadingComplexity\ of the full policy's \TotalComplexity\ mean effective
directions. Those averages describe annual histories, not the final-date
spectrum alone. Positive rank is calculated with a relative $10^{-12}$
tolerance, and near-tied eigenvalues are kept together at group boundaries.
The same operator also limits feasible complexity: ridge cannot activate
more directions than the managed matrix supports. This numerical rank
constraint is checked in the kernel and neural calculations. A long weak
tail can therefore change regularized exposure without constituting evidence
that its individual directions are reliably estimable or economically useful.

\begin{figure}[H]
\centering
\includegraphics[width=\textwidth]{figures/fig03_gaussian_complexity.pdf}
\caption{A large single-window illustration for the fixed Gaussian
representation. January 2024 historical and subsequent Sharpe are shown on
separate vertical scales, with effective complexity on the common horizontal
axis. The lower panel concatenates the 12 February 2024--January 2025 payoffs.
The star is a retrospective maximum over the frozen grid; the diamond is
selected using earlier validation $Q$. Neither estimates a significant
population optimum.}
\label{fig:complexity_gaussian}
\end{figure}

The final-year path isolates regularization: bandwidth, feature bank, and
historical sample remain fixed as $\lambda$ varies. Its retrospective peak
occurs at complexity \SinglePeakC\ with subsequent Sharpe \SinglePeakSR.
Historical fit and subsequent performance need not move together. A short
window can nevertheless yield a noisy shape, especially at aggressive
penalties. We therefore interpret this display as an illustration, then
examine all historical decision years in the next figure. No endpoint is
discarded, and a minimum-norm zero-penalty solution is retained separately
as an appendix diagnostic without changing validation eligibility.

Figure~\ref{fig:complexity_nine} supplies the historical context. Its first
eight blocks contain five decision years each; the final block contains seven.
For each penalty candidate we concatenate the monthly payoffs actually realized
across annual refits within the block, and calculate one Sharpe from that
series. We do not average annual Sharpes. Horizontal coordinates average the
corresponding annual complexities weighted by their realized month counts.
All annual windows here contain 12 months, so month and annual weighting
coincide. Curves use common axis definitions and terminate at their observed
supports.

All \InteriorCount\ of the nine grid maxima are interior, a descriptive
property of these curves rather than evidence that every period has a
statistically identified optimum. The maximizing mean complexity ranges
from \MinimumPeakC\ to \MaximumPeakC\, and the shapes and performance levels
vary markedly. The 2008--2012 peak is at relatively low complexity; later
windows have different supports and weaker realized performance. Broad
regions can be nearly flat relative to sampling uncertainty. The selected
annual policy is shown separately: its average complexity and realized Sharpe
need not lie on any fixed-$\lambda$ curve because its penalty varies annually.

Pointwise bands use 5,000 paired circular-block resamples, with the same
monthly draws across all candidates. They quantify uncertainty conditional
on realized paths under approximately stable dependence and distributional
conditions within a block. They do not correct for searching the grid or
establish validity under arbitrary historical nonstationarity. Six- and
24-month blocks, paired economic contrasts, and subperiod comparisons are
reported separately in the appendix.

As a secondary check, we jointly select bandwidth and ridge using prior
validation and multipliers $\{0.25,0.5,1,2,4\}$ of the initial median distance.
Each bandwidth's penalty scale comes from its own initial training spectrum.
The fixed control is numerically reproduced first. Gaussian fixed/tuned gross
Sharpes are \GaussianFixed\ / \GaussianTuned\ and net Sharpes are
\GaussianFixedNet\ / \GaussianTunedNet; the corresponding Mat\'ern values are
\MaternFixed\ / \MaternTuned\ and \MaternFixedNet\ / \MaternTunedNet.
Boundary selections and uncertainty remain part of this sensitivity, not a
reason to retune the main representation after observing test performance.

\clearpage
\begin{figure}[H]
\centering
\includegraphics[width=\textwidth]{figures/fig04_gaussian_nine_panels.pdf}
\caption{Nine non-overlapping Gaussian decision-year blocks, with fixed
bandwidth and feature bank. Panels report actual month counts. The first
block realizes February 1978--January 1983; each subsequent boundary follows
the same January-close convention, and the final block realizes February
2018--January 2025. Lines concatenate monthly fixed-$\lambda$ OOS payoffs
across refits. Stars mark retrospective grid maxima; diamonds show the
separate validation-selected annual policy. Shading gives pointwise 95\%
paired 12-month circular-block intervals, conditional on fitted paths and
uncorrected for peak selection. Full dates and candidate paths are public.}
\label{fig:complexity_nine}
\end{figure}

\clearpage
\subsection{What Additional Directions Trade}

We interpret portfolios associated with eigenspaces. At each annual refit,
eigendecomposition partitions the selected Gaussian ridge policy into four
disjoint positive-rank groups. Their actual spectral contributions become
stock-level signed holdings, which give weighted long and short characteristic
profiles and subsequent payoff and covariance attribution. The same selected
$\lambda$ is retained across nested comparisons. These are attribution
portfolios, not four separately optimized strategies, and equal-risk
diagnostic baskets do not support the economic claims.

\begin{figure}[H]
\centering
\includegraphics[width=\textwidth]{figures/fig05_economic_holdings.pdf}
\caption{Economic profiles of the actual Gaussian spectral holdings. The
left panel reports long-minus-short differences in weighted formation-rank
scores; the right reports signed notional exposures, preserving investment
scale. Families use an explicit signed characteristic dictionary. Debt
financing and equity issuance are distinct. Overlapping families describe
the same holdings and cannot be added as disjoint return components.}
\label{fig:economic_holdings}
\end{figure}

The profiles reveal reweighting within the same characteristic universe.
Leading holdings favor profitability and oppose recent returns and debt
financing. Intermediate directions modify those exposures rather than
providing an invariant named factor. A normalized characteristic contrast
can remain large when its component's notional is small; the two panels
therefore answer different questions. Formation-time descriptors characterize
positions, while realized outcomes determine their investment contribution.

\input{tables/spectral_value}

Table~\ref{tab:spectral_value} separates historical second-moment mass,
component gross, netted account gross, and subsequent investment performance.
The leading group has mean component gross \LeadingGross\ and annual excess
payoff \LeadingMean\%. Adding the intermediate group contributes
\IntermediateMean\% annually, yet its mean component gross \IntermediateGross\
raises cumulative account gross by only \IncrementalGross\. This is stock-level
cancellation: some new directions offset existing positions. The sum of
component gross exposures is consequently not the gross of the combined
account.

The added component has negative covariance with the leading portfolio.
Its annual covariance is \IntermediateCovariance\, but total annual variance
changes by \IntermediateVariance\. The account's variance can rise despite
negative covariance because the component adds its own variance. A positive
standalone mean similarly does not establish an improvement in portfolio
Sharpe; the combined mean, variance, and costs must be assessed together.

Here the gross Sharpe rises from \LeadingSR\ to \TwoGroupSR\, and the net
Sharpe from \LeadingNetSR\ to \TwoGroupNetSR. The paired 12-month-block
interval for the gross change is [\GroupGrossLow,\GroupGrossHigh], and for
the net change [\GroupNetLow,\GroupNetHigh]. These economic contrasts retain
their uncertainty and remain descriptive. They are not tests of a change in
$Q$, and no significant optimum or universal diversification benefit is
required for the attribution to be informative.

Small incremental account gross does not mean small turnover. Monthly
characteristics and annual refits can alter offsetting positions substantially.
Every nested account therefore sums group weights at the stock level before
calculating trades against drift-adjusted pretrade holdings, short borrowing,
and net returns. Separate component cost estimates would not reconstruct the
netted account. The appendix documents initial entry from cash, post-cost
target weights, cash financing, annual-to-monthly short-fee conversion, and
the upstream corporate-action treatment. Target-weight changes are reported
as a separate proxy and are never called realized trading costs.

The weaker tail has little historical second-moment mass and small absolute
payoff contributions. Its covariance can still matter for the selected
policy, but the financial evidence should be evaluated through cumulative
gross and net Sharpe rather than spectral mass alone. The same observation
limits interpretation of characteristic profiles: a striking normalized tilt
can arise in a direction whose investment contribution is economically small.
Alternative rank boundaries and feature-draw sensitivity are reported in the
appendix, without preserving labels when group content changes.

Six-benchmark regressions use the actual component returns. Full-sample
intercepts are descriptive spanning statistics, not implementable hedges or
newly identified risk premia. A separate hedge subtracts only slopes estimated
from each pretest history; it preserves the component's fitted scale. Detailed
coefficients and joint financing--profitability cells remain in the appendix.
Near-tied eigenvectors can rotate, so group profiles, rather than matched
individual eigenvector names, are the principal historical comparison.

The earlier history-length experiment remains a secondary diagnostic on its
common dates. It does not motivate a memory-selection policy. Economic
content and supported complexity vary with history. The stationary theory
benchmarks approximately stable conditions; a ``local'' interpretation requires
those conditions. Descriptive complexity changes do not identify changes in
$b$, and degrees of freedom do not certify profitable opportunities.
